% =============================================================
% Section 5: Discussion
% =============================================================
\section{Discussion}\label{sec:discussion}

\paragraph{Schwarzschild and Kerr field errors.}
The hFNO is our principal contribution, combining a
reduced-resolution numerical prior with correction operators trained
separately on the Schwarzschild and Kerr parameter sets. The coarse
finite-difference solve enforces the propagation, scattering and
boundary treatment of each perturbation equation, while the neural
operator corrects the discretisation error introduced by
under-resolution. Neither hFNO requires per-configuration
neural-network optimisation.

\paragraph{Evidence from Schwarzschild and Kerr.}
Across the Schwarzschild configurations, the hFNO reduces the mean
relative $L^2$ error from $13.39\%$ to $1.85\%$ and the field MSE by
a factor of $46$. Figure~\ref{fig:hyb-grad} shows a correlation
$r=0.921$ between $d_{\mathrm{idx}}$ and numerical-prior error, which
is concentrated near propagating wavefronts. For
Kerr, the median relative $L^2$ error falls from $59.5\%$ to $0.78\%$,
and the high-spin median frequency error is $0.38\%$. The
Richardson-extrapolated field reaches $0.008\%$, showing that the
remaining discrepancy arises primarily from the learned correction
rather than the target construction.

\paragraph{QNM extraction.}
Across the Schwarzschild configurations, M4 gives a median frequency
error of $0.152\%$. For Kerr, the retrospective Leaver-closest
selection gives median errors of $0.33\%$ in
frequency and $1.52\%$ in damping time, compared with $0.05\%$
and $0.27\%$ from the fine reference fields. Appendix~\ref{app:kerr-extractors}
shows the estimator dependence and distinguishes retrospective
selection from the pre-specified target-mode ensemble. The larger damping-time
error is consistent with training being dominated by the
high-amplitude burst and wavefronts, whereas damping depends on the
low-amplitude late-time envelope. The hFNO loss minimises full-field
mean-squared error and does not directly optimise QNM frequency or
damping time.

\paragraph{Scope and limitations.}
The test configurations lie within the parameter ranges used for
training, use one perturbative harmonic in each geometry and require
separate training for the Schwarzschild and Kerr formulations. The
fixed-grid Fourier representation has not been tested under
changes of domain, resolution or boundary treatment. The
gradient gate reduces the global Schwarzschild mean-squared error,
while the pointwise maps show smaller gains where the numerical prior
is already accurate.
Richardson extrapolation also relies on the two coarse
discretisations lying within the expected convergence regime.
The field results apply to the $\ell=2$ Zerilli and
$(\ell,m)=(4,2)$ Teukolsky sectors studied here; the reported QNM
errors concern their fundamental modes.

\paragraph{Future directions.}
The late-time errors in Figs.~\ref{fig:kerr-ringdown}
and~\ref{fig:kerr-pointwise} motivate late-time-sensitive objectives,
temporal decomposition or solver-in-the-loop training
\citep{Um2020solverloop,Wang2023longtime,McCabe2023stability}.
Initial data designed to excite stronger overtone content and
additional harmonics would test a broader range of ringdown signals.
End-to-end timing should measure the wall-clock reduction and training
break-even point.

Observer waveforms generated by the Schwarzschild and Kerr hFNOs
could also be incorporated into Bayesian inference. Reduced bases and
reduced-order quadratures
accelerate repeated gravitational-wave likelihood evaluations while
controlling the approximation error
\citep{Canizares2013roq,Canizares2015roq}. Constructing a reduced
basis from the hFNO waveforms would measure likelihood-evaluation cost
and propagate detector noise, hFNO field error and fit-window
dependence into parameter estimates.
